\section{Training Methodology}
\label{sec:training}

This section describes how the shared policy is actually trained: the
curriculum machinery that stages a run (Section~\ref{subsec:train-curriculum}),
the domain-randomization protocol (Section~\ref{subsec:train-dr}), the
exploration and learning-rate schedules (Section~\ref{subsec:train-schedules}),
a set of reinforcement-learning correctness decisions that the fault model
makes unusually delicate (Section~\ref{subsec:train-correctness}), and the
infrastructure that makes long runs reproducible and resumable
(Section~\ref{subsec:train-infra}).

\subsection{Curriculum}
\label{subsec:train-curriculum}

Training is organized as an ordered list of \emph{phases}, declared entirely
in the configuration file. Each phase specifies its episode count and may
override the team size, obstacle density, per-drone step budget, an
exploration reset, and --- crucially --- its own domain-randomization ranges;
anything not overridden falls back to the global settings. Two conventions
turn this list into a practical staging tool. First, \emph{a phase with zero
episodes is skipped}: alternative stagings are kept in the file and enabled
by editing a single number, so the exact shape of every run is versioned
alongside the code. Second, phase boundaries interact cleanly with resuming:
when a run restarts from a checkpoint, completed phases are skipped, the
active phase resumes from its saved episode index, and the phase's
exploration reset is suppressed (the checkpoint's annealed $\varepsilon$ is
kept rather than re-inflated).

The final run used in this report consists of a single main phase of
$50{,}000$ episodes with a per-drone budget of $200$ rounds and the full
randomization ranges of Table~\ref{tab:notation} active throughout ---
including the fault range --- preceded only by a zero-episode debug phase
kept for staging. Earlier exploratory stagings (shorter phases with narrower
ranges) informed the hyperparameters but do not contribute weights: the
reported policy is trained from scratch in the single phase. The curriculum
machinery is nevertheless part of the contribution: it is what made those
iterations cheap.

\subsection{Domain randomization}
\label{subsec:train-dr}

At the start of every episode the loop samples
$\theta = (\rvis, \rcomm, \nagents, \rho, \pfault)$ --- integers uniformly
(inclusive) and floats uniformly from their ranges --- applies it to the
environment via \texttt{set\_domain\_params}, and only then resets
(Section~\ref{subsec:arch-env}). The sampled ranges
(Table~\ref{tab:notation}) are wide by design. Vision spans myopic
($\rvis{=}1$: a $3\times3$ window) to far-sighted ($\rvis{=}6$); the
communication range spans nearly-isolated to easily-connected teams; the team
size spans a single drone --- where cooperation is moot and the policy must
search alone --- to four; obstacle density spans open fields to cluttered
mazes. The fault range $[0, 0.003]$ deserves its own reading: at the upper
end, a drone's probability of surviving all $200$ rounds is
$0.997^{200} \approx 0.55$, so the range spans ``faults never happen'' to
``losing roughly half the team over an episode is ordinary''. The policy is
therefore forced to be simultaneously a competent solo searcher, a
cooperative team member, and a survivor of team collapse --- with the
context vector (Section~\ref{subsec:arch-context}) telling it, at every
turn, which of these regimes it currently inhabits.

Per-episode resampling also acts as an implicit curriculum in expectation:
easy episodes (large vision, no faults, sparse obstacles) and hard ones
(blind, faulty, cluttered) interleave freely, which keeps a steady stream of
successful trajectories --- and hence of terminal reward signal --- flowing
into the replay buffer even before the policy masters the hard regimes.

\subsection{Schedules}
\label{subsec:train-schedules}

\paragraph{Exploration.}
$\varepsilon$-greedy exploration is annealed \emph{per episode} and
multiplicatively: starting from $\varepsilon_0 = 1$, each episode applies
$\varepsilon \leftarrow \max(0.05,\, 0.9995\,\varepsilon)$. The floor of
$0.05$ is reached after roughly $6{,}000$ episodes
($\approx 12\%$ of the run), after
which a residual $5\%$ of random actions keeps the buffer from collapsing
onto the greedy policy's own trajectory distribution. Exploration coins are
drawn per drone per turn (Section~\ref{subsec:arch-learning}), so within one
round some drones may explore while the others act greedily.

\paragraph{Learning rate.}
The learning rate follows a linear warmup ($10\%$ of the run,
$10^{-4} \to 5\times10^{-4}$) into a cosine annealing down to $10^{-6}$ at
the final episode. Two details are deliberate. The schedule is stepped
\emph{once per episode} rather than per gradient step: episode lengths vary
wildly across randomized regimes (a lucky find ends an episode in tens of
turns; a blind, faulty maze consumes the full budget), so stepping per
update would tie the schedule's shape to the difficulty mix rather than to
training progress. And it is sized by the \emph{total} episode count across
all phases, so the decay is a property of the run, not of any single phase.

\paragraph{Update cadence.}
One Double-DQN gradient step (Section~\ref{subsec:arch-learning}) runs every
$4$ agent-turns on minibatches of $16$ transitions from the shared replay
buffer; the target network is synchronized by hard copy every $500$ gradient
steps. Appendix~\ref{app:config} consolidates all values.

\subsection{Reinforcement-learning correctness decisions}
\label{subsec:train-correctness}

The fault model interacts with value-based learning in ways that are easy to
get silently wrong: nothing crashes, training even improves, but the learned
values answer a subtly different question than intended. This subsection
records the decisions taken, each with its rationale; together they
constitute the project's learning-side treatment of faults.

\paragraph{Time-outs bootstrap.}
Transitions record \emph{termination} only; running out of budget is a
truncation, and the final window still bootstraps the next-state value
(Section~\ref{subsec:bg-dqn}; \citealp{pardo2018time}). The time limit is an
experimental artifact --- the world does not end at turn $T_{\max}$ --- and
treating it as terminal would teach the network that states near the budget
boundary have no future, biasing values along every long trajectory.

\paragraph{A fault is a truncation, not a terminal.}
When a drone breaks, its pending $n$-step window is flushed immediately
\emph{with} bootstrapping ($\gamma^k$ for the $k$ steps accumulated), and no
further transitions are generated for it --- wrecks' no-op turns never enter
the buffer. The rationale mirrors the time-out case but is easier to get
wrong, because a fault \emph{feels} terminal for that drone. It is not, in
the sense that matters for learning: the fault strikes independently of the
drone's state and action ($\pfault$ is exogenous), so the states that
happened to precede a fault are not worth less than identical states that
did not. Zeroing the bootstrap would charge the value function for bad luck,
systematically depressing the values of \emph{whatever the drone happened to
be doing} when it broke --- noise injected as signal, growing with
$\pfault$. Bootstrapping instead treats the fault as censoring: the return
is simply observed no further. The alternative --- modelling faults
\emph{inside} the value function by scaling values by survival probability
--- would require the network to know $\pfault$, which is exactly the
quantity Section~\ref{subsec:arch-context} argues a deployed drone cannot
know.

\paragraph{Team credit is retroactive, and only for the living.}
With the shared terminal reward enabled
(Section~\ref{subsec:prob-objective}), the finder's terminal reward arrives
through the environment, while every \emph{other} live drone's bonus is
written retroactively into its most recent still-pending window, which is
then marked terminal. This mechanism has two consequences worth making
explicit. It requires $n \geq 2$: with one-step windows every transition is
emitted immediately, and there is nothing pending to credit (the
configuration is validated at run start). And it excludes drones that broke
before the find --- their windows were flushed at fault time, so their
experience of the episode ends, correctly, with what they themselves
observed: a drone that died exploring an empty quadrant should not learn
that dying there is worth the team's eventual success.

\paragraph{Selection is protected from shaping drift.}
Every $1{,}000$ episodes the current policy is evaluated near-greedily
($\varepsilon = 0.05$) on a \emph{fixed} set of $20$ instances: the
randomization draw and the map layout are both derived deterministically
from fixed seeds, and one seed reproduces layout and fault sequence alike
(Section~\ref{subsec:arch-env}), so successive evaluations differ only in
the policy. The checkpoint-selection rule is success-dominant: a new best
requires a strictly higher success rate, with mean reward used as a
tie-breaker only at $100\%$ success. This shields model selection from the
shaped reward terms of Table~\ref{tab:reward} --- a policy must not be
preferred because it harvests exploration bonuses more efficiently while
finding the target less often.

\paragraph{No planner contamination.}
The BFS auto-navigation override of Section~\ref{subsec:arch-tooling} never
runs during training --- neither in the episode loop nor in the in-training
evaluation that selects checkpoints. In the loop it would corrupt the data
distribution (transitions would reflect the planner's policy, not the
network's); in the selection it would corrupt the metric (checkpoints would
be ranked partly on the planner's competence). It exists strictly as an
analysis instrument for Section~\ref{sec:experiments}.

\subsection{Infrastructure and reproducibility}
\label{subsec:train-infra}

Training runs on a SLURM-managed HPC cluster (single GPU per run), launched
by a batch script that stages temporary files on node-local scratch storage;
the identical entry point runs unchanged on a workstation, which is how
short debugging stagings were iterated. Full runs log per-episode scalars to
TensorBoard --- reward, length, success, the sampled randomization values,
the number of drones lost to faults, $\varepsilon$ and the learning rate ---
so a run's difficulty mix and its learning progress can be inspected jointly.

\paragraph{Checkpointing and resumption.}
Three kinds of checkpoint are written: \texttt{best.pt} (weights only, on
every new evaluation best), periodic per-phase snapshots, and
\texttt{last.pt}, which additionally carries the full training state ---
optimizer, learning-rate schedule position, $\varepsilon$, gradient-step
counter, and curriculum position (phase and episode index). Resuming from
\texttt{last.pt} therefore continues the run \emph{as if uninterrupted},
with one deliberate exception: the replay buffer is not serialized (at
$2\times10^5$ full observation pairs it would dwarf the weights), so a
resumed run refills its buffer from fresh experience. A separate
weights-only warm-start mode initializes a brand-new run (fresh optimizer,
schedules, and curriculum) from an existing policy. On checkpoint
compatibility: the observation layout and context dimensionality changed
when the fault model was introduced (new Broken channel, new belief entry),
so weights from the pre-fault ancestor of this project are architecturally
incompatible --- training for this report starts from scratch.

\paragraph{Parallel checkpoint evaluation.}
Because the in-training evaluation is intentionally small (20 instances), a
separate \emph{watcher} process --- spawned automatically at run start ---
evaluates every periodic checkpoint on $200$ randomized-but-seeded episodes
as soon as it is written, appending to an incremental, crash-safe summary
and producing a final ranking when training ends. This offloads thorough
evaluation from the training loop (which would otherwise stall for minutes
every $1{,}000$ episodes) and, in practice, provides a second opinion on
checkpoint selection at no wall-clock cost to training.

% Figure TODO (after the final run): training curves — episode reward,
% eval success rate, epsilon/LR schedules, n_broken histogram.
